% Section 2: Reinforcement learning for forecasting
\subsection{Reinforcement learning for forecasting}
\label{sec:rl_for_forecasting}

\subsubsection{Reinforcement learning as an ensemble combiner}
\label{sec:rl_ensemble}

The simplest way to put reinforcement learning near forecasting is to treat the
dynamic combination of an ensemble of base forecasters as a sequential decision
problem. \citet{Fu2022} formalise this as a Markov decision process and call the
resulting agent RLMC. Given $N$ pretrained base forecasters
$\mathcal M = \{M_1, \dots, M_N\}$, the state at time $t$ is
$s_t \in \mathbb R^{T \times d_s}$, namely the recent $T$-step window of the time
series together with a log of recent per-model performance. The action
$a_t = (w^1_t, \dots, w^N_t) \in \Delta^{N-1}$ is a vector of non-negative weights
on the simplex, and the realised forecast at time $t$ is the convex combination
$\hat y_t = \sum_{i=1}^N w^i_t \hat y^i_t$ of the base predictions. The transition
satisfies
\begin{equation}
P(s_{t+1} \mid s_t, a_t) = P(s_{t+1} \mid s_t),
\label{eq:fu2022_transition}
\end{equation}
that is, the agent's weighting choice does not move the underlying time series.
The reward is a mixture of a sMAPE-based and a rank-based normalised score,
$r_t = \alpha R^{\mathrm{sMAPE}}_t + (1-\alpha) R^{\mathrm{rank}}_t$, where each
term is renormalised to $[-1, 1]$ by quantile transformation against the
in-sample distribution of base-model errors at time $t$.

\citet{Fu2022} draw two methodologically useful observations from this setup. The
first concerns exploration. Because the action does not enter the state-transition
kernel by~\eqref{eq:fu2022_transition}, the conditional next-state distribution can
be sampled from the empirical state-trajectory of the training data with any
action superimposed. The combination problem is therefore in effect a model-based
RL problem in which the transition kernel and the reward are both known up to the
empirical sample, and an unbounded supply of transitions
$(s_t, a_t, r(s_t, a_t), s_{t+1})$ can be generated without further interaction. The
second observation concerns the structure of the optimal policy. If only one base
model is best at each state, the optimal policy is nearly deterministic, namely
$\pi^\star(s) \approx e_{i^\star(s)}$, the standard basis vector at the best model
index. A policy-gradient method then drives $\log \pi_\phi(a \mid s) \to -\infty$
on most of the action space, and the estimator
$\hat g = \mathbb E_t[\nabla_\phi \log \pi_\phi(a_t \mid s_t) \hat A_t]$ becomes
numerically unstable. Figure 2 of \citet{Fu2022} documents the divergence
empirically on the ETT dataset.

These observations motivate the choice of Deep Deterministic Policy Gradient
\citep{LillicrapDDPG2016} as the optimisation algorithm. The actor
$a = \pi_\phi(s)$ outputs a continuous deterministic action, and the critic
$Q_\theta(s, a)$ approximates the optimal action-value function under the
mean-squared Bellman loss
$\mathcal L_Q(\theta) = \mathbb E_{(s,a,r,s')}[ (Q_\theta(s,a) - r - \gamma
Q_{\theta'}(s', \pi_\phi(s')))^2 ]$, with the policy updated by deterministic
policy gradient
$\nabla_\phi J = \mathbb E_s[ \nabla_a Q_\theta(s, a)\rvert_{a = \pi_\phi(s)}
\nabla_\phi \pi_\phi(s) ]$. The deterministic actor sidesteps the log-probability
explosion, and the off-policy structure exploits the model-based generation route
in the first observation. Three exploration tricks complete the algorithm. The
actor is pretrained on a multi-class classification surrogate whose labels are
the indices of the per-state best base model. An $\epsilon$-greedy rule with a
convex-combination inductive bias replaces the standard Gaussian perturbation,
namely $a = \pi_\phi(s)$ with probability $1 - \epsilon$ and $a = e_i$ for a
uniformly drawn $i$ with probability $\epsilon$. A second replay buffer holds
hard transitions whose realised reward falls below a threshold $\hat r_t$, and
mini-batches are drawn from both buffers so that the easy distribution does not
absorb the agent toward a degenerate copy-the-best-model policy.

The empirical evidence in \citet{Fu2022} covers four datasets. On the ETT
electricity-transformer-temperature series, RLMC reports a sMAPE of $72.11$ against
$73.69$ for the single best base model and $79.83$ for the uniform average. On
the M4 daily panel, RLMC reports a sMAPE of $2.31$ against $2.39$ for the single
best and $3.93$ for uniform. On the Climate series the MAE falls to $1.47$ against
$1.67$ for single best, and on the GEFCOM electricity-load benchmark RLMC ties
with FFORMA at MAE $88$ to $87$. The gains are real but modest, in the range of
one to three per cent of sMAPE or MAE relative to the single best base model,
and are sometimes overtaken on individual datasets by the gradient-boosted
FFORMA baseline. The ablation in Figure 7 of \citet{Fu2022} suggests that the
exploration strategy carries most of the improvement, with the pretraining and
the hard-sample buffer contributing marginal additional gains.

The methodological depth of this strand is modest, and we mention it primarily to
demarcate the scope of the chapter. The model combination problem in
\citet{Fu2022} can be recast as an online learning problem on the simplex with a
non-stationary loss sequence. The regret guarantees one would obtain from a Hedge
or an EXP3 algorithm against the best fixed combination in hindsight are
$O(\sqrt{T \log N})$ and $O(\sqrt{N T \log N})$ respectively, and there is no
indication that DDPG-based weight learning improves on these rates. The genuine
contribution of \citet{Fu2022} is the demonstration that off-policy actor-critic
machinery can be applied to model-combination problems with imbalanced
best-model distributions, not a regret bound. The deeper RL-for-forecasting
strands, namely Brier-loss reward training, approximate dynamic programming on a
sequential-estimation MDP, and decision-focused learning, sit on richer
methodological foundations and occupy the remainder of this section.

\subsubsection{Reinforcement learning with decision-relevant rewards}
\label{sec:rl_brier}

A second strand of RL-for-forecasting uses on-policy reinforcement learning to fine
tune a probabilistic forecaster against a proper scoring rule that doubles as the
downstream decision loss. The cleanest production-scale instance to date is
\citet{Turtel2025}, who adapt the family of reinforcement-learning-with-verifiable-rewards
algorithms, which had previously been confined to deterministic mathematical
verification, to the stochastic and lagged-supervision setting of real-world
forecasting. The base model is the 14B-parameter DeepSeek-R1-Distill-Qwen reasoning
model, the training corpus consists of $10{,}000$ historical Polymarket yes/no
questions augmented with $100{,}000$ synthetic questions generated by the Lightning
Rod Foresight framework, and the evaluation is a held-out set of $1{,}265$ test
questions with strict precautions against temporal leakage.

The forecasting episode is a single decision. The agent receives a textual prompt
that includes the question and a set of news headlines dated strictly before the
prediction date, and emits a single probability $\hat p \in [0, 1]$ for the event.
After the resolution date, the binary outcome $y \in \{0, 1\}$ is revealed and the
episode terminates with reward $R = -(\hat p - y)^2$. The reward is the negative
Brier score of \citet{Brier1950}, a strictly proper scoring rule in the sense of
\citet{GneitingRaftery2007}, so the population-optimal policy is the true
conditional probability $\hat p = \mathbb P(y = 1 \mid \text{prompt})$. The training
schedule is strictly online and single-pass, with questions encountered in
chronological order and outcomes revealed at the resolution date. Multiple epochs
are not used because re-exposure to a resolved question lets the model recover the
future from training-time information.

The algorithmic contribution is to identify why standard GRPO destabilises in this
regime and how to fix it. Group Relative Policy Optimisation \citep{Shao2024deepseek}
draws $G$ rollouts $\{o_1, \dots, o_G\}$ from the old policy and forms the advantage
$\hat A_i = (r_i - \mu) / \sigma$, where $\mu$ and $\sigma$ are the mean and
standard deviation of the rewards across the group. The per-question $\sigma$
normalisation is helpful when rewards are scale-heterogeneous across questions, but
\citet{Turtel2025} argue that it is harmful when the reward is a Brier score because
the asymmetry between modest correct-and-overconfident gains and rare large
miscalibration penalties is precisely the signal the model needs to learn proper
calibration. Their Modified-GRPO drops the $\sigma$ division and uses the
baseline-subtracted advantage $\hat A_i = r_i - \mu$, paralleling the Dr.\,GRPO
modification of \citet{Liu2025DrGRPO}. ReMax \citep{LiReMax2023} replaces the group
mean with a learned baseline $b_i$ and uses $\hat A_i = r_i - b_i$, which removes
the variance normalisation while preserving an unbiased advantage estimate. A
malformed model output that fails the regex extraction is assigned the maximum
Brier loss of $1.0$ during training (strict Brier) and a constant $0.25$ during
evaluation (soft Brier), with the latter chosen to preserve gradient signal without
inducing training collapse.

The failure modes that motivated these modifications are documented qualitatively
and quantitatively. Standard GRPO without guardrails drives the model toward
extreme overconfidence, with $39.3\%$ of probabilities falling into the $0$-$10\%$
or $90$-$100\%$ buckets and many landing on exact zero or one. Modified-GRPO
without guardrails reduces the extreme fraction to $7.9\%$ but introduces a
different pathology, namely code-switching between Chinese and English mid-sequence
and ungrammatical partial outputs. Adding the strict-Brier reward, a token-length
cap, an early-stopping gibberish filter, and the ReMax baseline brings the extreme
fraction down to $13.1\%$ while restoring grammatical coherence. The optimisation
scaffold across algorithms is AdamW with the standard reasoning-LLM hyperparameters
and bfloat16 precision, with global gradient-norm clipping at $1.0$ and an entropy
bonus coefficient of $0.001$. Total compute is roughly three days on eight NVIDIA
H100 GPUs per training run.

The headline empirical results vindicate the design. On the $1{,}265$-question hold
out, the seven-run ReMax ensemble trained on the $110{,}000$-question combined
corpus attains a soft-Brier score of $0.190$ with $95\%$ confidence interval
$[0.178, 0.203]$ and an expected calibration error of $0.062$, against $0.202$ and
$0.093$ for OpenAI o1, $0.215$ and $0.089$ for the base DeepSeek-R1 14B model, and
$0.151$ and $0.043$ for the contemporaneous Polymarket market price. The Brier
improvement over o1 is $-0.011$ with two-sided $p < 0.05$, and the calibration
improvement is large in absolute terms but does not reach statistical significance
against Polymarket itself. None of the trained models match the market on Brier,
which is consistent with the market aggregating information across many traders,
but the trained forecasters retain incremental information beyond the price.

The decision-side evaluation establishes the practical relevance. \citet{Turtel2025}
construct a hypothetical trading rule that places a one-share bet on every test
question and measure cumulative profit on Polymarket. The ReMax ensemble earns
$\$52$ in simulated profit on $\$433$ of committed capital, corresponding to
roughly a $12\%$ return on investment. The Modified-GRPO model earns $\$54$ on
$\$428$, and OpenAI o1 earns $\$39$. The profit concentrates on questions where
the market was most uncertain. On the subset with market probability between $40\%$
and $60\%$, the ReMax ensemble's bets win $11.8$ percentage points more often than
the market price predicts, with $95\%$ confidence interval $[8.0, 15.7]$ and a
return on investment of roughly $20\%$. The pattern is the empirical realisation of
the Granger-Pesaran argument from
Section~\ref{sec:forecasting_rl_intro}. A forecaster trained to minimise a proper
scoring rule, with the reward instrument set to the same Brier loss that the
evaluator computes, recovers calibration even when the prediction-accuracy
performance does not surpass an aggregating market, and the calibration translates
into realised profit on the threshold-crossing decisions where the market is
indifferent.

The methodological lesson generalises beyond LLMs and prediction markets. Treating
a strictly proper scoring rule as the reward signal of a sequential RL agent
converts the calibration question into a policy-optimisation question, which is
amenable to standard variance-reduction techniques. The two failure modes
identified by \citet{Turtel2025}, namely extreme overconfidence under raw GRPO and
gibberish under raw Modified-GRPO, are not specific to LLMs but to the interaction
of variance-normalised policy gradients with bounded reward distributions. A
forecasting workflow that emits continuous quantile predictions and trains under
the continuous ranked probability score of \citet{GneitingRaftery2007}, or that
emits a categorical distribution and trains under the log score, can be expected
to encounter the same gradient instabilities and to benefit from the same
baseline-subtraction modifications. The connection between Brier-loss RL training
and the SPO programme of Section~\ref{sec:dfl} is direct. Both are special cases
of the differentiable-Wald-decision-theory recipe, with the difference being that
Brier-loss RL trains under the scoring rule itself rather than under the realised
decision regret of a specific downstream optimisation.

\subsubsection{Approximate dynamic programming for forecasting}
\label{sec:adp_forecasting}

Sequential estimation has a longer pedigree in the dynamic programming tradition.
\citet{BertsekasRollout2022} provides the unifying frame. A measurement policy at
time $t$ is a function from the current posterior over an unknown $\theta$ to the
next observation point, and rollout improves a base policy by performing a one-step
lookahead under Monte Carlo simulation of the residual horizon. The certainty
equivalence approximation replaces the stochastic continuation cost with its
expectation conditional on the realised first step, which lowers the computational
burden without forfeiting the policy improvement guarantee whenever the base policy
already satisfies a mild contractivity condition. The framework covers Bayesian
optimisation, optimal stopping for sequential testing, and adaptive Bayesian
experiment design, and it returns the same vocabulary as standard RL.

\citet{Fernandes2024} applies the neuro-dynamic programming variant of this idea to
the gross domestic product nowcasting problem. The reward is not the cross-sectional
forecast error. It is the adjustment cost between successive predictions plus a
terminal accuracy penalty, which Fernandes argues better reflects the way a central
bank or a finance ministry uses a nowcast. The value function is approximated by a
single-hidden-layer feedforward network with $32$ units, trained by temporal
difference learning on quarterly data from $26$ European Union countries with
Portugal held out for evaluation. Over the period 2015Q1 to 2023Q4, the neuro-dynamic
programming nowcaster lowers the out-of-sample mean absolute error from $6.6$ to
$4.5$ percentage points of GDP range, a $32\%$ reduction over OLS, with the root
mean squared error falling from $8.3$ to $5.9$ percentage points of range, a $28\%$
reduction. The most striking finding is that training on foreign data transfers well
to Portugal, which alleviates the small-T problem that plagues macroeconomic
forecasting.

\citet{BenAmeur2025} extend the dynamic programming side of the literature into the
online and non-stationary regime. They give the first sublinear regret bounds for
online decision-focused learning. The DF-FTPL algorithm achieves static regret of
order $T^{-1/4}$ against the best fixed parameter in hindsight, and the DF-OGD
algorithm achieves dynamic regret of order $(1 + P_T)^{1/4} T^{-1/4}$, where $P_T$
is the path variation of the comparator. Both bounds are independent of the
parameter dimension up to a $\log \log d$ factor, which matters for high-dimensional
forecasting models.

\subsubsection{Decision-focused learning}
\label{sec:dfl}

The deepest sub-stream is decision-focused learning, which trains the forecast
model under a loss that reflects the downstream decision regret rather than the
prediction error. The framework was formalised by \citet{ElmachtoubGrigas2022} under
the name Smart Predict-then-Optimize, abbreviated SPO. Given features $x$ and a
linear cost vector $c$ that determines the objective $c^\top z$ of a downstream
optimisation problem $\min_{z \in \mathcal Z} c^\top z$, the SPO loss measures the
suboptimality of the decision induced by the predicted cost $\hat c(x)$, namely
$\mathrm{SPO}(\hat c, c) = c^\top z^\star(\hat c) - c^\top z^\star(c)$, where
$z^\star(c)$ is the optimal solution under the true cost. The SPO loss is
discontinuous in $\hat c$, so \citet{ElmachtoubGrigas2022} construct the convex
surrogate $\mathrm{SPO}^+(\hat c, c)$. Their main theorem establishes Fisher
consistency.

\begin{theorem}[\citealp{ElmachtoubGrigas2022}, Theorem 2]
\label{thm:spo_fisher}
Suppose the noise around the conditional mean cost $\mathbb E[c \mid x]$ is centrally
symmetric and the feasible region $\mathcal Z$ is a convex polyhedron. Then the
population minimiser of the SPO$^+$ loss is the conditional mean cost function. In
particular, the regression function $x \mapsto \mathbb E[c \mid x]$ minimises both
the Bayes risk of the SPO loss and the Bayes risk of the SPO$^+$ surrogate.
\end{theorem}

The Fisher consistency result is non-trivial, because the SPO$^+$ loss is not a
proper scoring rule in the sense of \citet{GneitingRaftery2007}. Symmetric noise is
needed to align the surrogate's minimiser with the SPO minimiser, and asymmetric
noise generically breaks the equivalence.

\citet{LiuGrigas2021} convert Fisher consistency into a finite-sample statement. They
prove that on convex polyhedra the SPO$^+$ calibration function $\delta(\varepsilon)$
satisfies $\delta(\varepsilon) \ge c \varepsilon^2$, which combined with a standard
Rademacher complexity bound on the surrogate excess risk yields an excess SPO risk
of order $n^{-1/4}$ for empirical risk minimisation over a Rademacher-bounded
hypothesis class. For strongly convex level sets the rate improves to $n^{-1/2}$,
matching the calibration of ordinary squared error.

\citet{DontiAmosKolter2017} operationalised the SPO programme for neural networks by
differentiating through the Karush-Kuhn-Tucker conditions of the downstream
optimisation. They apply task-based learning to three economic problems, namely
electricity grid scheduling on PJM data, energy storage arbitrage, and an inventory
stock balancing problem, and report regret reductions of $10\%$ to $40\%$ over models
trained with standard maximum likelihood and plugged into the same optimisation
routine. The implicit differentiation step is the methodological bridge from
batch DFL to deep learning at scale, and it is the building block on which most
subsequent DFL papers rest.

The DFL story is not uniformly positive. \citet{HuKallus2020} prove a counterintuitive
result for contextual linear optimisation under low near-dual-degeneracy. Under their
margin condition, the regret of a naive plug-in estimator that fits the conditional
mean by ordinary least squares decays at rate $n^{-1}$, while an integrated decision-aware
estimator achieves $n^{-2/3}$. The naive method exploits the curvature of the loss
near the decision boundary in a way that a decision-aware method does not. Their
result is a useful corrective. Decision-focus is a worthwhile asymptotic objective
only when the hypothesis class cannot recover the true conditional cost, which is
the realistic regime for economic forecasting but is not the same as the universal
claim that DFL dominates.

The modern frontier addresses misspecification head on. \citet{HuangGupta2024}
introduce the family of Perturbation Gradient losses. The PG loss is a Lipschitz
continuous difference-of-concave surrogate whose stochastic gradient approximates
the directional derivative of the plug-in decision loss via zeroth-order
estimation. The key theorem states that under a sub-Gaussian noise model and a
Rademacher-complexity bound $R_n$ on the hypothesis class, the population
minimiser of the PG loss attains a regret of order
$\widetilde O\bigl((dp/n)^{1/4}\bigr)$ relative to the best-in-class policy. The
guarantee is best-in-class rather than full-information optimal, but it holds under
misspecification, which is precisely the case where SPO$^+$ loses its Fisher
consistency property. Empirically, PG outperforms SPO$^+$ on synthetic and
real-data instances where the cost function does not lie in the hypothesis class.

\citet{Mandi2024} survey the broader landscape. They organise gradient-based
methods, comprising surrogate-loss methods like SPO$^+$ and implicit-differentiation
methods like that of \citet{DontiAmosKolter2017}, alongside gradient-free methods such
as smoothing and Bayesian approaches. They benchmark $11$ DFL techniques on $7$
canonical problems, including shortest path, knapsack, portfolio, travelling
salesman, and bipartite matching, and report that DFL is most useful precisely when
the underlying prediction problem is misspecified. On well-specified problems,
the gap between predict-then-optimise and decision-focused learning shrinks to
within statistical error.
